% Generated from locale/tr/content/lambda-calculus/syntax/alpha.tex; do not hand-edit.


\olfileid[tr]{lam}{syn}{alp}
\olsection{$\alpha$-Dönüşümü}

$\lambd[x][x]$ ile $\lambd[y][y]$ arasında nasıl bir ilişki vardır? Her
ikisi de özdeşlik fonksiyonunu gösterir. Elbette sözdizimsel olarak farklı
terimlerdir. Yalnızca bağlı değişkenin adı bakımından ayrılırlar; biri,
diğerindeki bağlı değişkenin ``yeniden adlandırılmasıyla'' elde edilir. Buna
\emph{$\alpha$-dönüşümü} denir.

\begin{defn}[Bağlı değişkeni değiştirme, $\aconvone$]
  Bir $M$ terimi $\lambd[x][N]$ geçişini içeriyor, $y\notin\FV{N}$ ve
  $\Subst{N}{y}{x}$ tanımlıysa, bu geçişin
  \begin{equation*}
    \lambd[y][\Subst{N}{y}{x}]
  \end{equation*}
  ile değiştirilerek $M'$ teriminin elde edilmesine \emph{bağlı değişkeni
  değiştirme} denir ve $M\redone[\alpha]M'$ yazılır.
\end{defn}

\begin{defn}[Bir bağıntının bağdaşması]
  Terimler üzerindeki bir $R$ bağıntısı aşağıdaki koşulları sağlıyorsa
  \emph{bağdaşır} denir:
  \begin{enumerate}
  \item $R N N'$ ise $R \lambd[x][N] \lambd[x][N']$;
  \item $R P P'$ ise $R (PQ) (P'Q)$;
  \item $R Q Q'$ ise $R (PQ) (PQ')$.
  \end{enumerate}
\end{defn}

Tanımı şimdi başka biçimde ifade edelim:
\begin{defn}[Bağlı değişkeni değiştirme, $\aconvone$]
  \emph{Bağlı değişkeni değiştirme} ($\redone[\alpha]$), terimler üzerinde
  aşağıdaki koşulu sağlayan en küçük bağdaşır bağıntıdır:
  \begin{align*}
    &\lambd[x][N] \redone[\alpha] \lambd[y][\Subst{N}{y}{x}] && \text{$x
      \neq y$, $y \notin \FV{N}$ ise} \\
    & &&\text{ve $\Subst{N}{y}{x}$ tanımlıysa.}
  \end{align*}
\end{defn}

Buradaki ``en küçük'', bağıntının yalnızca bağdaşma ve ek koşulun
gerektirdiği çiftleri içerdiği, başka hiçbir çift içermediği anlamına gelir.
Dolayısıyla bağıntı şöyle de tanımlanabilir:

\begin{defn}[Bağlı değişkeni değiştirme, $\aconvone$] \ollabel{defn:aconvone}
  \emph{Bağlı değişkeni değiştirme} ($\aconvone$) tümevarımlı olarak şöyle
  tanımlanır:
  \begin{enumerate}
  \item $N \aconvone N'$ ise $\lambd[x][N] \aconvone
    \lambd[x][N']$. \ollabel{defn:aconvone1}
  \item $P \aconvone P'$ ise $(PQ) \aconvone (P'Q)$.
    \ollabel{defn:aconvone2}
  \item $Q \aconvone Q'$ ise $(PQ) \aconvone (PQ')$.
    \ollabel{defn:aconvone3}
  \item $x \neq y$, $y \notin \FV{N}$ ve $\Subst{N}{y}{x}$ tanımlıysa
    $\lambd[x][N] \redone[\alpha] \lambd[y][\Subst{N}{y}{x}]$.
    \ollabel{defn:aconvone4}
  \end{enumerate}
\end{defn}

Tanımlar eşdeğerdir; ispatı alıştırma olarak bırakıyoruz. Bundan sonra
tümevarımlı tanımı kullanacağız.

\begin{defn}[$\alpha$-dönüşümü, $\aconv$]
  \emph{$\alpha$-dönüşümü} ($\aconv$), $\aconvone$ bağıntısını içeren,
  terimler üzerindeki en küçük yansımalı ve geçişli bağıntıdır.
\end{defn}

Yukarıdaki gibi ``en küçük'', bağıntının yalnızca geçişliliğin ve
$\aconvone$ bağıntısının gerektirdiği çiftleri içerdiği anlamına gelir. Bu da
aşağıdaki eşdeğer tanıma götürür:

\begin{defn}[$\alpha$-dönüşümü, $\aconv$] \ollabel{defn:aconv}
  \emph{$\alpha$-dönüşümü} ($\aconv$) tümevarımlı olarak şöyle tanımlanır:
  \begin{enumerate}
  \item $P \aconv Q$ ve $Q \aconv R$ ise $P \aconv R$.
    \ollabel{defn:aconv1}
  \item $P \aconvone Q$ ise $P \aconv Q$. \ollabel{defn:aconv2}
  \item $P \aconv P$. \ollabel{defn:aconv3}
  \end{enumerate}
\end{defn}

\begin{ex}
  $\lambd[x][f x]$, $\lambd[y][f y]$ terimine $\alpha$-dönüşür; tersi de
  geçerlidir. Gayriresmî olarak ikisi de bir argüman alıp bu argümanın $f$
  altındaki değerini döndüren ve ortam değişkeni $f$'ye gönderme yapan
  fonksiyonlardır.

  $\lambd[x][f x]$, $\lambd[x][g x]$ terimine $\alpha$-dönüşmez.
  Gayriresmî olarak bunlar sırasıyla $f$ ve $g$ ortam değişkenlerine gönderme
  yapar; $f$ ile $g$'nin farklı olduğu ortamlarda farklı davrandıkları için
  farklı fonksiyonlardır.
\end{ex}

\begin{prob}
  Aşağıdaki terim çiftleri $\alpha$-dönüşümlü müdür?
  \begin{enumerate}
  \item $\lambd[x][\lambd[y][x]]$ ve $\lambd[y][\lambd[x][y]]$;
  \item $\lambd[x][\lambd[y][x]]$ ve $\lambd[c][\lambd[b][a]]$;
  \item $\lambd[x][\lambd[y][x]]$ ve $\lambd[c][\lambd[b][a]]$.
  \end{enumerate}
\end{prob}

\begin{lem}\ollabel{lem:fv-one}
  $P \aconvone Q$ ise $\FV{P} = \FV{Q}$.
\end{lem}

\begin{proof}
  $P \aconvone Q$ !!{derivation} üzerinde tümevarımla.
  \begin{enumerate}
  \item Son kural \olref{defn:aconvone4} ise $P$, $\lambd[x][N]$;
    $Q$ ise $\lambd[y][\Subst{N}{y}{x}]$ biçimindedir; burada $x\neq y$,
    $y\notin\FV{N}$ ve $\Subst{N}{y}{x}$ tanımlıdır. $x\in\FV{N}$ olup
    olmamasına göre durumları ayırırız:
    \begin{enumerate}
    \item $x \in \FV{N}$ ise:
      \begin{align*}
        \FV{\lambd[y][\Subst{N}{y}{x}]} &=
        \FV{\Subst{N}{y}{x}} \setminus \{y\} \\
        & = ((\FV{N} \setminus \{x\}) \cup \{y\}) \setminus \{y\}
         && \text{\olref[sub]{thm:infv} gereği} \\
        & = \FV{N} \setminus \{x\} \\
        & = \FV{\lambd[x][N]}.
      \end{align*}
    \item $x \notin \FV{N}$ ise:
      \begin{align*}
        \FV{\lambd[y][\Subst{N}{y}{x}]}
        & = \FV{\Subst{N}{y}{x}} \setminus \{y\} \\
        & = \FV{N} \setminus \{y\}
         && \text{\olref[sub]{thm:notinfv} gereği} \\
        & = \FV{N} \\
        & = \FV{N}\setminus\{x\} \\
        & = \FV{\lambd[x][N]}.
      \end{align*}
    \end{enumerate}
  \item Diğer üç durum alıştırma olarak bırakılmıştır.
  \end{enumerate}
\end{proof}

\begin{prob}
  \olref[lam][syn][alp]{lem:fv-one} lemasının ispatını tamamlayın.
\end{prob}

\begin{lem}\ollabel{lem:inv}
  $P \aconvone Q$ ise $Q \aconvone P$.
\end{lem}

\begin{proof}
  $P \aconvone Q$ !!{derivation} üzerinde tümevarımla.
  \begin{enumerate}
  \item Son kural \olref{defn:aconvone4} ise $P$, $\lambd[x][N]$ ve $Q$,
    $\lambd[y][\Subst{N}{y}{x}]$ biçimindedir; burada $x\neq y$,
    $y\notin\FV{N}$ ve $\Subst{N}{y}{x}$ tanımlıdır. Önce
    \olref[sub]{thm:clr} gereği
    $y\notin\FV{\Subst{N}{y}{x}}$ olur. \olref[sub]{thm:inv} gereği
    $\Subst{\Subst{N}{y}{x}}{x}{y}$ yalnızca tanımlı değil, aynı zamanda
    $N$'ye eşittir. Dolayısıyla \olref{defn:aconvone4} gereği
    $\lambd[y][\Subst{N}{y}{x}]
    \aconvone \lambd[x][\Subst{\Subst{N}{y}{x}}{x}{y}] =
    \lambd[x][N]$ elde edilir.
  \end{enumerate}
\end{proof}

\begin{prob}
  \olref[lam][syn][alp]{lem:inv} lemasının ispatını tamamlayın.
\end{prob}

\begin{thm}
  $\alpha$-dönüşümü terimler üzerinde bir denklik bağıntısıdır; yani
  yansımalı, simetrik ve geçişlidir.
\end{thm}

\begin{proof}
  \begin{enumerate}
  \item Her $M$ terimi, bağlı değişkenlerde \emph{sıfır} değişiklik yapılarak
    $M$'ye dönüştürülebilir.
  \item $P$, bir dizi bağlı değişken değişikliğiyle $Q$'ya
    $\alpha$-dönüşüyorsa, bu değişiklikleri \olref{lem:inv} gereği ters
    yönde ve ters sırada uygulayarak $Q$'dan $P$ elde edilir.
  \item $P$, bir değişiklik dizisiyle $Q$'ya; $Q$ da başka bir değişiklik
    dizisiyle $R$'ye $\alpha$-dönüşüyorsa, önce ilk, sonra ikinci diziyi
    uygulayarak $P$'yi $R$'ye dönüştürebiliriz.
  \end{enumerate}
\end{proof}

Bundan sonra $M$, $N$'ye $\alpha$-dönüşüyorsa, $M$ ile $N$'nin
\emph{$\alpha$-denk} olduğunu, $M\aeq N$, söyleriz. Az önce gösterdiğimiz
gibi bu, ancak ve ancak $N$ de $M$'ye $\alpha$-dönüşüyorsa geçerlidir.

\begin{thm}\ollabel{thm:fv}
  $M \aeq N$ ise $\FV{M} = \FV{N}$.
\end{thm}

\begin{proof}
  Doğrudan \olref{lem:fv-one} lemasından çıkar.
\end{proof}

\begin{lem}\ollabel{lem:sub:R}
  $R \aeq R'$ ve $\Subst{M}{R}{y}$ tanımlıysa $\Subst{M}{R'}{y}$ de
  tanımlıdır ve $\Subst{M}{R}{y}$ ile $\alpha$-denktir.
\end{lem}

\begin{proof}
  Alıştırma.
\end{proof}

\begin{prob}
  \olref[lam][syn][alp]{lem:sub:R} lemasını ispatlayın.
\end{prob}

\olref[sub]{sec} içinde yerine koymanın bazı durumlarda tanımsız olduğunu
hatırlayın. Bununla birlikte terimler üzerinde $\alpha$-dönüşümü kullanarak
bağlı değişkenleri yeniden adlandırabilir ve yerine koymayı her zaman tanımlı
hâle getirebiliriz. Aşağıdaki teorem sonucun $\alpha$-denkliğini koruduğunu
gösterir:

\begin{thm}\ollabel{thm:sub}
  Her $M$, $R$ ve~$y$ için $M \aeq M'$ ve $\Subst{M'}{R}{y}$ tanımlı
  olacak biçimde bir $M'$ vardır. Ayrıca $M''\aeq M$, $R''\aeq R$ ve
  $\Subst{M''}{R''}{y}$ tanımlı olacak biçimde başka bir $M'',R''$ çifti
  varsa $\Subst{M'}{R}{y} \aeq \Subst{M''}{R''}{y}$ olur.
\end{thm}

\begin{proof}
  $M$'nin kuruluşu üzerinde tümevarımla:
  \begin{enumerate}
  \item $M$, bir $z$ değişkenidir: alıştırma.
  \item $M$'nin $\lambd[x][N]$ biçiminde olduğunu varsayın. $x$ ve $y$'den
    farklı, ayrıca $z\notin\FV{N}$ ve $z\notin\FV{R}$ olacak biçimde bir
    $z$ değişkeni seçin. Tümevarım varsayımına göre $N'\aeq N$ ve
    $\Subst{N'}{z}{x}$ tanımlı olacak biçimde bir $N'$ vardır.
    \olref{defn:aconvone}\olref{defn:aconvone1} gereği
    $\lambd[x][N]\aeq\lambd[x][N']$ olur. Ayrıca
    \olref{defn:aconvone}\olref{defn:aconvone4} gereği
    $\lambd[x][N']\aeq\lambd[z][\Subst{N'}{z}{x}]$ olur. Bunu yapabiliriz;
    çünkü $z\neq x$, $N'\aeq N$ olduğundan $z\notin\FV{N'}$ ve
    $\Subst{N'}{z}{x}$ tanımlıdır. Son olarak
    $\Subst{\lambd[z][\Subst{N'}{z}{x}]}{R}{y}$ tanımlıdır; çünkü
    $z\neq y$ ve $z\notin\FV{R}$.

    Aynı koşulları sağlayan başka bir $N''$ ve $R''$ varsa
    \begin{multline*}
      \Subst{(\lambd[z][\Subst{N''}{z}{x}])}{R''}{y} \aeq\\
    \begin{aligned}
      &\lambd[z][\Subst{\Subst{N''}{z}{x}}{R''}{y}] \\
      &\aeq \lambd[z][\Subst{\Subst{N''}{z}{x}}{R}{y}] && \text{
        \olref{lem:sub:R} gereği}\\
      &\aeq\lambd[z][\Subst{\Subst{N'}{z}{x}}{R}{y}]
      && \text{tümevarım varsayımı gereği}\\
      &\aeq\Subst{(\lambd[z][\Subst{N'}{z}{x}])}{R}{y}.
    \end{aligned}
    \end{multline*}
    \item $M$, $(PQ)$ biçimindedir: alıştırma.
  \end{enumerate}
\end{proof}

\begin{prob}
  \olref[lam][syn][alp]{thm:sub} teoreminin ispatını tamamlayın.
\end{prob}

\begin{cor}\ollabel{cor:sub}
  Her $M$, $R$ ve~$y$ için $M'\aeq M$, $R\aeq R'$ ve
  $\Subst{M'}{R'}{y}$ tanımlı olacak biçimde bir $M',R'$ çifti vardır.
  Ayrıca $M''\aeq M$, $R''\aeq R$ ve $\Subst{M''}{R''}{y}$ tanımlı olacak
  biçimde başka bir çift varsa
  $\Subst{M'}{R'}{y}\aeq\Subst{M''}{R''}{y}$ olur.
\end{cor}

\begin{proof}
  Doğrudan \olref{thm:sub} teoreminden çıkar.
\end{proof}

